\section{Discussion}
\label{sec:discussion}

% Remaster, 3 October 2026.

\subsection{What the results say about harness design}

Compared as shipped, the six-agent AutoDS-Tools ranks first on TabReD and
the workflow FEDOT.LLM and the minimal harness Terminus-2 are level. The
ablations attribute that order to two things that are separable from the
harness: the instruction file of one system and the backend search of
another. With the file removed, the three designs, from the fixed-stage
workflow to the open shell, reach the same region of the normalised scale,
and the differences between them equal the range a single harness covers
over its own attempts. Three designs cannot fix the shape of a curve, so we
do not claim that a designed middle ground would beat all three. What the
data support is narrower: on tabular tasks with a mid-size open model, the
choice of harness changes accuracy by less than the choice of attempt, and
changes reliability, wall-clock and the set of tasks a system can attempt
at all.

The instruction file is what changes accuracy, and the ablations show its
limits. It pays inside the harness it was written for, from
scratch, on the task family whose library it names. Moved to another
harness it ends half or more of the trials; applied where a working script
already exists it ends whole tasks. Both failures trace to specific sentences, the
advice on training time and the instruction to replace the supplied tool,
and both are invisible to an ablation that switches the file as one block.
This is also the reading we propose for the disagreement among published
ablations of instruction files and retrieved knowledge: a file measured in
one harness, from one starting point, has not been measured in another.

What separates every configuration from the strongest published method is
search. The agents choose the right model family unaided and write one
configuration; the reference methods are Optuna-tuned over fifteen seeds.
The one harness that searches, FEDOT.LLM, beats every published method on
the task where model choice matters most and is level with Terminus-2 over
the eight. Search over models and hyperparameters is a known, cheap addition
to any of the three harnesses, and the traces say it is the first thing to
add.

\subsection{Guidance for practitioners}

The results give one rule per situation. When a strong tuned baseline
exists, use it: the best cell in the paper is level with tuned XGBoost and
every cell is below the strongest published method. For a tabular task from
scratch with accuracy first and an hour per task acceptable, use
AutoDS-Tools with its instruction file. For a first pass over a new
scientific dataset in minutes per attempt, use Terminus-2 with the task text
alone, which beat two author baselines and found the published
featurisation on OpenPoly. Where search over a fixed backend pays, use
FEDOT.LLM, best of the three on \texttt{ecom-offers} and on F-DATA. When a
working script is supplied, add no library instruction. Two rules hold
across at least two protocols: prescribe only the decision the model would
otherwise get wrong, and keep advice about training time out of the prompt
of any harness with a hard time limit.

\subsection{Limitations}
\label{sec:threats}

\emph{Repeats.} Published baselines are means over fifteen seeds; agent rows
are means of three attempts, and on \allIdentFedot{} of eight tasks the three
FEDOT.LLM attempts agree to every digit, so its spread is understated. With
eight tasks the paired tests have little power, and we state them alongside
the win counts and intervals for that reason.

\emph{Scorers that accept degenerate submissions.} The prohibition on
constant submissions belongs to the instruction file, so only the cell with
the file carried it, and any bias from the asymmetry works against that
cell. Both degenerate submissions we found belong to Terminus-2: on
\texttt{imdb} one label for all 25\,000 test rows, scored at exactly
$0.5000$ and excluded, and on \texttt{amp-parkinsons} one global value in
$61.5\%$ of rows, read as a partial fallback.

\emph{The repair loop.} AutoDS-Tools repairs failed executions more often
with the file than without (Section~\ref{sec:friction}), so it partially
masks the cost of the file. We report the counts without adjustment;
disabling the loop in both branches is the experiment we would run next.

\emph{The lost trial.} One of the 24 AutoDS-Tools trials on TabReD
reached the one-hour limit with nothing to score and is excluded, so
\texttt{homesite-insurance} rests on two attempts. Scoring that attempt at
the weakest published value would move the normalised mean of AutoDS-Tools
from $\nAutoDS$ to $\sensLostNorm$, still ahead of the other two
harnesses; the exclusion drops the slowest attempt of a system whose
failures fall on the tasks it finds hardest.

\emph{Data access.} Several MLAgentBench tasks use subsets of the original
data and none an official test split, so every comparison there is
internal. On TabReD the agents receive the labelled validation split
(Section~\ref{sec:benchmarks}), a small advantage that works against our
main conclusion.

\emph{The backbone.} The backbone ablation rests on three models, two from
one vendor, and one task family. A stronger model that reasons less
verbosely might not push the multi-agent system into the time limit at all,
and we cannot separate capability from verbosity with what we ran. The
finding that the instruction file outweighs a backbone swap is therefore
reported for this domain and these models.

\emph{Generality of the harnesses.} The three systems are one instance each
of a workflow, a multi-agent system and a minimal harness. A different
six-agent system or a different fixed-stage workflow could behave
differently; the claim is about these three, and about the method of
separating harness, instructions and backbone, which transfers.
